\section{Engineering and verification}\label{sec:engineering}
Part~6 holds the rules about how the platform is built and shown to be correct: source provenance, code quality, dependencies, module structure, tests, performance and release gates. These rules apply from the first commit (phase P0). The first deliverable is a usable object-storage and fine-grained authorization service with a client SDK that encrypts before upload and decrypts only on the authorized client. Platform servers process application-content ciphertext and opaque recipient envelopes; they never receive content plaintext, plaintext data-encryption keys or recipient private keys. MCP, model submission and external tool execution consume its contracts in later milestones; their packages and deployments are not prerequisites for the storage foundation.

\subsection{Independent implementation and source research}\label{sec:source-research}
The platform is an independently designed Go object-storage and access-control system. MinIO and Cerbos are research references: their mapped behavior informs design decisions, while the project authors its own implementation, tests, contracts and documentation. Neither upstream is a production component. The aim is to retain useful mechanisms without inheriting either product's code base, deployment model or full feature surface.

\req{SRC-001}{P0}{REV}{Implementation code, tests and expressive documentation are independently authored. MinIO and Cerbos source, libraries, generated implementation files, tests, comments and assets \mustnot\ be copied, translated, vendored, linked, embedded or packaged into this product. Their servers \mustnot\ be used as production runtimes behind a wrapper, and this project is not a fork or rebranding. Renaming symbols, changing language or machine-assisted paraphrasing does not satisfy this requirement. Studying behavior does not authorize incorporation of its expression.}
\req{SRC-002}{P0}{REV}{For each adopted or adapted mechanism, a research record identifies the upstream repository and revision, the smallest decisive inspected source or documentation range, observed behavior and limits, and the mapping feature ID. A separate original design record defines the project's contract, algorithm or state transitions, failure behavior, owning module and acceptance invariant, and explains what is retained, changed or rejected. The mechanism adoption matrix links these records to the existing verification harness. Source snapshots and copied research excerpts stay outside production source and release artifacts.}
\req{SRC-003}{P0}{REV}{Maintained general-purpose libraries may be admitted through \rr{LIB-002} and \rr{LIB-003}; this permission does not admit MinIO or Cerbos implementation packages or product runtimes. Each admitted dependency has a targeted provenance and license record, pinned version and retained required notices. The distribution license for original project code is chosen separately from obligations attached to dependencies or other included material. Independent authorship and this research process do not automatically establish license compatibility or remove third-party rights.}
\reqn{SRC-004}{P0}{REV}{Before affected implementation work starts, a qualified reviewer records the license and provenance constraints at each studied or admitted revision and any required separation of research and implementation work. No upstream executable is used as a runtime, test oracle or reference (D-01). Original requirements and independently authored test inputs are the acceptance authority.}

\subsubsection{From mapped behavior to original Go code}
The working sequence is: inspect one decisive mapping entry; record the mechanism in behavioral terms; design the project's smaller contract; implement that contract independently; and verify its security and failure invariant in the shared harness. Reopen upstream evidence only when a concrete ambiguity blocks the design. Do not read several equivalent sources or reproduce an upstream subsystem merely to imitate its structure.

\textbf{Inspiration, not transcription.} Ideas, behaviors and invariants observed upstream may be reused; upstream code, structure and expression may not (\rr{SRC-001}). Code is written from this specification and the project\textquotesingle s records, never from upstream files or from recollection of them. A mechanism taken from the studied projects is recorded before it is coded: the repository\textquotesingle s mechanism matrix (\texttt{docs/research/mechanism-matrix.md}) is the \rr{SRC-002} record for the mappings made so far, with one decision per row (adopt, adapt, reject or defer) and the owning contract. A mechanism that has no row gets one first.

The reviewable deliverables are a pinned mechanism note, an original design decision, the owning module and contract, and the smallest distinct acceptance case. A maintainer checks both semantic fidelity and implementation independence. Existing cases are extended when they already exercise the invariant; the research matrix is not a request for a new test per row. Unselected mechanisms remain deferred, rather than becoming unused packages or configuration switches.

\subsubsection{Selected mechanism adoption matrix}
Table~\ref{tab:mechanisms} is the adoption baseline for the selected mechanisms. It summarizes the existing functionality map, not a new full-source audit or an assertion of measured upstream behavior. The complete map retains file coverage and unresolved observations. Changes to an adoption decision update both the original design record and this table.

The evidence set is \texttt{map\_cerbos.json}, \texttt{map\_minio\_storage.json} and \texttt{map\_minio\_access.json} in the separate research report. Cerbos is pinned to \code{ae014ba2ea919771216bc666c4f654d18a1c1a52}; MinIO is pinned to \code{7aac2a2c5b7c882e68c1ce017d8256be2feea27f}. Each source link below points to a decisive anchor already recorded by that map. The anchors justify the observed mechanism, not the project's independently chosen design.

\begin{longtable}{@{}P{0.21\textwidth}P{0.39\textwidth}P{0.10\textwidth}P{0.20\textwidth}@{}}
\caption{Selected upstream ideas and their original implementation}\label{tab:mechanisms}\\
\toprule Observed mechanism & Our design and owner & Choice & Invariant and evidence\\\midrule\endfirsthead
\toprule Observed mechanism & Our design and owner & Choice & Invariant and evidence\\\midrule\endhead
\href{https://github.com/cerbos/cerbos/blob/ae014ba2ea919771216bc666c4f654d18a1c1a52/internal/ruletable/check.go#L429}{Cerbos: principal, role and scope evaluation paths} & Our own evaluator in \code{policy/} and \code{decision/} for \code{fga.policy/v1}: a matching deny overrides an allow across roles and scopes, a per-principal rule is a condition, and a separate always-run platform ceiling cannot be bypassed. Scope ancestry does not establish tenant identity. & Reject upstream precedence; adapt scope ancestry & An extra role never undoes a deny, and a ceiling objection is never lifted. T-01, T-02.\\
\href{https://github.com/cerbos/cerbos/blob/ae014ba2ea919771216bc666c4f654d18a1c1a52/internal/ruletable/manager.go#L105}{Cerbos: validate, build and replace policy state} & \code{bundle/} builds an immutable candidate off the request path, validates it, and activates one signed epoch. Failed activation does not make the candidate ready; old state cannot serve requests requiring a newer epoch. & Adapt & Every request and batch uses one authorized snapshot; failed builds and stale nodes cannot imply success. T-14.\\
\href{https://github.com/cerbos/cerbos/blob/ae014ba2ea919771216bc666c4f654d18a1c1a52/internal/ruletable/index/index.go#L56}{Cerbos: narrow rule candidates and reuse expression programs} & \code{policy/} indexes only necessary dimensions and uses a qualified general-purpose CEL library. Begin with a clear reference evaluation path; add indexing only where measurements justify it. No upstream index or cache implementation is reused. & Adapt & An optimization cannot omit a relevant rule or change an effect. Compare paths in T-01; bound costs in T-02/T-22.\\
\href{https://github.com/minio/minio/blob/7aac2a2c5b7c882e68c1ce017d8256be2feea27f/cmd/object-api-interface.go#L246}{MinIO: separate object operations from physical storage} & \code{content/}, \code{publication/} and \code{backend/} expose small typed contracts. PostgreSQL owns logical visibility; immutable backend objects become readable only after an authorized catalog commit. MCP later calls this same service. & Adapt & A backend write alone grants no access. Conditional publication binds the exact version and current authority. T-06, T-08.\\
\href{https://github.com/minio/minio/blob/7aac2a2c5b7c882e68c1ce017d8256be2feea27f/cmd/erasure-object.go#L1440}{MinIO: quorum durability and corruption detection} & R1 uses our own single-node store with the flush discipline of \rr{BKD-001}. R3 \code{replicated/} uses three full replicas and two durable domain receipts, with exact digests and placement bindings. Erasure coding and MinIO disk formats are not adopted. Checksums do not substitute for authorization or authenticated encryption. & Adapt; defer erasure coding & No acknowledged publication without the required durability evidence; no automatic quorum reduction. T-15, T-16, T-17.\\
\href{https://github.com/minio/minio/blob/7aac2a2c5b7c882e68c1ce017d8256be2feea27f/cmd/erasure-server-pool-decom.go#L605}{MinIO: resumable repair and evacuation} & R3 \code{maintenance/} persists bounded work in PostgreSQL. Copy, verify, change placement by compare-and-swap, then retire. A restricted workload receives a fenced maintenance plan; it does not impersonate a content reader. & Adapt & Restarts cannot lose required work, and stale workers cannot retire live copies. T-17, T-20.\\
\href{https://github.com/minio/minio/blob/7aac2a2c5b7c882e68c1ce017d8256be2feea27f/cmd/admin-handlers-users.go#L2893}{MinIO: derived credentials with parent restrictions} & \code{identity/} and \code{delegation/} bind actor, represented subject, tenant, audience and expiry. Delegation can narrow authority only. One native policy authority serves every storage and service route; separate IAM engines and owner bypasses are not adopted. & Adapt & A service identity or child credential cannot expand the represented principal's authority. T-03, T-20.\\
\href{https://github.com/minio/minio/blob/7aac2a2c5b7c882e68c1ce017d8256be2feea27f/cmd/encryption-v1.go#L554}{MinIO: envelope keys and streaming object protection} & Server mode (D-12): the storage role encrypts each version under a fresh data key sealed by a tenant KEK that only the \code{keys} role holds and unseals only under a signed decision verified in the authority fence; the key is bound to tenant, object, version and blob rather than a path. Client mode: \code{sdk/go/crypto/} encrypts content and constructs recipient-bound HPKE envelopes, and the \code{keys} role authorizes delivery of existing opaque envelopes only. No MinIO encryption code or wire format is reused. & Adapt & Server-mode keys are released only under exact fresh decisions; server credentials cannot decrypt client-mode content; clients reject wrong tenant, unit or recipient bindings. T-10, T-11, T-12, T-23.\\
\bottomrule
\end{longtable}

\subsubsection{Research evidence and its limits}
The mechanism matrix retains the pinned source anchors and the behavior each informed. The earlier targeted inspection covered six short excerpts (158 lines); its detailed file/line ledger is retained in the source bundle as \texttt{engineering-notes.tex}. That inspection supports the stated behavioral observations, not runtime durability, performance, compatibility or security qualification. Implementation structure is chosen from the project\textquotesingle s own contracts and measured needs.

\subsection{Code quality and naming}
\req{QUAL-001}{P0}{REV}{Security and correctness are release priorities ahead of source-code size, delivery speed or convenience. Maintainable complexity is optimized, not source compression. Dead code, speculative abstractions, duplicate implementations, unused configuration and unnecessary tests are removed, while default-deny enforcement, mandatory encryption, validation, error handling, bounded resources, durability checks and audit evidence are kept even when they cost code. No hard function-length, total-size or coverage-percentage quota overrides review judgment.}
\req{QUAL-002}{P0}{REV}{Code exposes coherent domain responsibilities and keeps distinct security concepts in distinct types, including actor versus represented subject and logical versus cryptographic version. Contracts use typed requests, explicit states and validated outcomes. Names make authority and side effects clear. Package layout and individual function names are implementation choices provided ownership, reviewability and these boundaries remain clear.}
\req{QUAL-003}{P0}{REV}{Trust assumptions, preconditions, side effects, ownership and lifetime, cancellation and failure outcomes are documented where they matter. \code{AuthorizeOperation} does not perform the protected write, and \code{Prepare} does not dispatch. A transactional function identifies its commit point and unknown-outcome behavior. \code{context.Context} is passed through I/O and blocking work, and the owner of each stream and buffer is defined. Errors carry stable classifications, preserve causes for internal diagnostics, and are redacted for clients. A panic is never recovered into an allow, and an irreversible effect is never retried without its idempotency contract.}
\req{QUAL-004}{P0}{REV}{Changes to authentication, policy semantics, delegation, key handling, publication, replication, fencing or external side effects receive independent technical review by someone qualified for the area. AI-assisted code receives the same review, and compilation is not evidence of a correct security design. CI runs formatting, \code{go vet}, scoped static analysis, relevant unit and integration scenarios, the race detector on meaningful concurrent paths, vulnerability scanning and targeted fuzzing of parsers, signatures, envelopes and cursors. Unresolved findings are recorded and an unmitigated applicable security or correctness defect blocks release.}
\req{QUAL-005}{P0}{REV}{A rushed or hard-to-explain patch is not accepted to meet a date. Review may require simplification, a better-supported dependency, clearer names or more evidence, and a change to security behavior needs explicit requirements and migration notes, not an incidental refactor.}

\subsection{Libraries and dependencies}
R1 content encryption and decryption run in the client SDK. Object bytes remain ciphertext during upload, storage, replication, backup and download, including after TLS terminates at a platform server. Infrastructure encryption additionally protects permitted catalog metadata and backups, but it does not replace client encryption. The envelope broker never unwraps a data key. Client private keys remain under the recipient or organization's control; an optional client-controlled KMS or HSM integration requires its own qualification and is not a server dependency. R1-C extends the same design to independently encrypted units, and R2 adds MCP over the ciphertext contract.

The libraries in Table~\ref{tab:libs} are selection candidates, not completed security reviews. Exact versions, checksums and toolchain compatibility are recorded at admission in \texttt{dependencies.lock.json}.
\begin{longtable}{@{}P{0.24\textwidth}P{0.32\textwidth}P{0.34\textwidth}@{}}
\caption{Candidate libraries and what the project still owns}\label{tab:libs}\\
\toprule Concern & Candidate & Project responsibility\\\midrule\endfirsthead
\toprule Concern & Candidate & Project responsibility\\\midrule\endhead
HTTP, TLS, context, logs & Go standard library & Listener, time and resource bounds, workload identity, redaction, role isolation\\
Policy expressions & CEL for Go & Environment, schema and error modes, platform obligations, immutable releases\\
Client streaming and recipient encryption & Tink for Go in the client SDK & Authenticated object/unit format, recipient binding, client key custody and recovery\\
Catalog and pooling & pgx; sqlc as a build-time generator & Explicit SQL transactions, tenant predicates, compare-and-swap and fencing\\
OIDC and JOSE & go-oidc; go-jose & Issuer, audience, rotation, expiry, delegation and freshness checks\\
Schema validation, YAML & A maintained JSON Schema library and Go YAML implementation & Pinned dialect, bounded decoding, duplicate and unknown-field rejection, no network schema loading\\
MCP transport (R2) & Official MCP Go SDK & Ciphertext/resource-envelope handlers; authorized client SDK or local bridge performs decryption\\
Internal RPC & Protobuf with qualified Go generators & One canonical contract, bounded messages, mTLS, deadlines\\
Catalog HA (R1) & PostgreSQL with an established HA coordinator & Verified quorum and fencing, backups, restore\\
S3 client conformance (R3) & AWS SDK for Go v2 & Independent protocol verification; the client signer is not a server verifier\\
\bottomrule
\end{longtable}
\req{LIB-001}{P0}{REV}{A standard or third-party component must fit the required semantics. Go's \code{ServeMux} cleans and redirects some paths, so S3 signing and object-key semantics need a raw handler. Schema validation after decoding cannot recover a duplicate JSON field already discarded. Library defaults are tested or configured explicitly, not assumed secure.}
\req{LIB-002}{P0}{REV}{Every direct dependency has an owner and a short record covering why it is needed; alternatives and the standard library; maintainers and support policy; release and security-response history; current advisories; test and conformance evidence; API stability; license and notices; transitive footprint; and the upgrade or replacement path. Popularity or frequent commits do not demonstrate maintenance quality, and a stable library is not rejected merely because it changes infrequently. Narrow maintained libraries are preferred over a large framework used for one helper.}
\req{LIB-003}{P0}{REV}{\code{go.mod}, \code{go.sum}, code generators and deployment images are pinned, an SBOM is produced and required notices are preserved. Vulnerability analysis runs in CI and periodically after release, and an applicable vulnerability requires a fix, removal, isolation or a documented time-bounded disposition. Updates are reviewed and tested and then deployed with the staged rollout and rollback discipline of product changes. Automatic unreviewed production upgrades are avoided, a necessary security update is not delayed to preserve a pin, and each dependency has security-notification routing and a tested replacement plan.}

\subsection{Module and role structure}
The implementation keeps the following ownership boundaries. The file and package layout is illustrative; a different decomposition is acceptable when it preserves the same authority, deployment and dependency boundaries.
\begin{longtable}{@{}P{0.24\textwidth}P{0.66\textwidth}@{}}
\toprule Owner & Responsibility\\\midrule\endhead
Decision and control & Verified identity and delegation, policy semantics and release state, grants, approvals and audit authority\\
Storage & Canonical content operations, publication and catalog consistency, private backend access, document and lineage metadata, server-mode payload encryption and decryption\\
Keys role & Tenant KEKs and decision-bound data-key wrap and release (server mode); recipient public-key records and authorized opaque-envelope publication and delivery (client mode); no blobs or payload plaintext\\
Client SDK & Authenticated content transfer; for client mode, trusted key custody, payload encryption/decryption, origin verification and client-side ingestion\\
Later adapters and services & MCP and S3 translate protocols; optional execution and replication preserve the same resource and authority contracts\\
\bottomrule
\end{longtable}
Cohesive code stays together; packages split when a real boundary warrants it. Future services do not require empty packages or deployment scaffolding in R1.
\req{MOD-001}{P0}{T-20}{The product uses one Go module with explicit role entrypoints and shared contracts. It includes no MinIO or Cerbos implementation dependency or production runtime. R1 provides independently deployed control, storage, decision and \code{keys} services plus the client SDK, which encrypts for client mode; R2 adds MCP and, only when selected, external execution; R3 may add custom blob storage. Separate entrypoints compose shared services without duplicating policy or content semantics. Policy validation and authorized administration remain available operational interfaces.}
\req{MOD-002}{P0}{T-20}{Each production process selects one role and initializes only its declared routes, dependencies, credentials and background work. Storage, MCP, envelope delivery, optional tool execution and blob serving run under separate identities, mounts, network policy and replica sets. Recipient key custody and client-mode content decryption run exclusively on authorized clients, outside every server role; server-mode payload encryption and decryption run only in the storage role, and KEK custody only in the \code{keys} role. MCP can scale without adding storage capacity, and storage can scale without deploying MCP. An all-roles mode, if provided, is development-only and rejected by production configuration, and packaging grants neither MCP nor tools direct backend access or recipient private keys. No other server role has data-key unwrap or content-decryption credentials, and no role can apply them to client-mode content.}
\req{MOD-003}{P0}{T-20}{Production configuration is typed, versioned and validated before listeners become ready, with deterministic file, environment and flag precedence. Unknown fields, insecure missing values and incompatible role capabilities are rejected. Decision and replica RPC require authenticated workload identity. Playground, reflection, debug and profiling endpoints, payload logging and administrative listeners are disabled unless explicitly needed, isolated and authorized. Reloadable settings are enumerated, and security-authority, recipient-registry and listener-role changes require a controlled restart or a separately reviewed activation path.}
\reqn{MOD-004}{P0}{T-20}{Each role is built from an explicit entrypoint and has a reviewed transitive dependency allowlist. Storage, control and decision do not depend on MCP, model-provider clients or external-tool execution. Every server binary excludes client-mode payload-decryption and HPKE unwrap code, including the envelope broker. Server-mode payload encryption and decryption code is linked only into the storage role and KEK wrap and unwrap code only into the \code{keys} role; other server cryptography is limited to its declared transport, signatures, ciphertext digests, public-key validation and opaque-format checks. CI verifies these dependency boundaries, and changed permissions or imports require review. Roles roll out independently with compatible contracts.}

\subsection{Keeping the code base small}
\req{SMALL-002}{P0}{REV}{Repetitive API clients and DTOs are generated from one canonical definition per contract, and SQL access uses reviewed SQL plus generated query bindings rather than hand-maintained parallel data-access models. Tool and resource schemas and wire contracts remain distinct domains. Generators are pinned, output is reproducible and generated files are not manually patched.}
\req{SMALL-003}{P0}{REV}{R1 starts with our own single-node blob store and the \code{local} profile, one client cryptographic profile and an opaque-envelope delivery service. R2 introduces MCP over the existing ciphertext contract and one authorized client SDK or local bridge. A client-controlled KMS or HSM is an optional integration, not an initial server service; one real downstream connector is required only if the optional external-tool service is selected. Extend through narrow provider interfaces when requirements justify it. No plugin framework, custom cryptography, replacement database, custom consensus or broad protocol implementation is justified solely by potential future use.}
\req{SMALL-004}{P0}{REV}{Handwritten production code, generated code, tests, direct and transitive dependencies and duplicated security logic are tracked. New subsystems are reviewed for maintenance cost, and dead code and unused dependencies are removed. Maintainability and attack surface are measured as well as lines of code, because moving complexity into unchecked generated code or a large dependency is not automatically a reduction.}
\reqn{SMALL-005}{P0}{REV}{PostgreSQL provides the catalog, publication intents, required approvals and the audit outbox in R1; later features reuse it for reservations, invocation ownership and bounded repair queues while its measured capacity suffices. Future-feature tables and workers are added when the feature is implemented, not as empty scaffolding. Redis, Kafka or a workflow engine is added only for a concrete requirement with a reviewed failure model. SQL leases, idempotency and outbox delivery still need explicit implementation and tests, since a library does not supply the project's transaction semantics.}

\subsection{Testing}\label{sec:testing}
\req{TEST-001}{P0}{REV}{Every retained test demonstrates a distinct required behavior, security invariant, failure mode or regression. Table~\ref{tab:tests} is an evidence checklist and not an instruction to create a test per row, function, requirement, adapter or layer. Near-duplicate tests are consolidated into table-driven scenarios, trivial accessors, generated code and library internals are not tested, and external behavior is preferred over mocks and snapshots. Test code meets the same naming, ownership and review standards as production code, and no coverage percentage or test count is a product objective.}
\req{TEST-002}{P0}{REV}{Similar failure cases are retained only when they protect distinct guarantees and the distinction is stated. Wrong-tenant access, an expired delegation, a stale commit generation, a lost quorum and ciphertext truncation exercise different boundaries; they may share a table-driven harness but none is proven by the success path alone. A regression test reproduces the defect with the smallest stable case, and fuzz targets focus on hostile input boundaries.}
\reqn{TEST-003}{P0}{LINT}{Each acceptance case declares its minimum phase, required features and status in \texttt{vectors/}. A \code{scenario_only} entry records an obligation but is not executable conformance evidence. A \code{fixture_ready} entry binds concrete initial state, inputs, ordered steps and expected assertions; the implementation harness consumes those inputs without inventing case-specific facts. Fault cases first establish their successful prerequisites. The checker validates identifiers, applicability and fixture structure. Release qualification additionally requires applicable fixtures and passing execution records bound to the current fixture, selected immutable release manifest and complete enabled configuration, plus evidence for non-vector tests. The manifest identifies the qualified role binaries; records from different releases or configurations cannot be combined. A structural check cannot establish implementation conformance. Changing normative inputs or expectations changes the specification version.}
\reqn{TEST-004}{R1-A}{T-01}{Policy-evaluation vectors run against the evaluator in CI. A change to the evaluator's semantics changes the vectors and the specification version together.}

\subsubsection{Security and correctness matrix}\label{sec:testmatrix}
\begin{longtable}{@{}P{0.06\textwidth}P{0.09\textwidth}P{0.68\textwidth}P{0.09\textwidth}@{}}
\caption{Acceptance tests}\label{tab:tests}\\
\toprule Test & Phase & Required evidence & Vectors\\\midrule\endfirsthead
\toprule Test & Phase & Required evidence & Vectors\\\midrule\endhead
\testrow{T-01}{R1-A, R1-B}{One conformance corpus covers deny-overrides across roles and scopes (an extra role never undoes a deny), per-principal rules as conditions, platform ceilings (a business ALLOW cannot override a ceiling DENY) and default denial. A ceiling with no matching deny yields neutral PASS: a business ALLOW remains allowed and a business DENY remains denied. Indexed and straightforward execution agree. Requests through different adapters, including the native object API in R1-B, produce identical canonical decisions; public responses disclose only authorized metadata, including on denial.}{yes}
\testrow{T-02}{R1-A}{Strict evaluation is the only mode: a bundle that selects a lenient mode or any construct the format does not define is rejected at bundle validation; a ceiling set containing an allow rule, or a manifest that omits its ceiling declaration, is rejected; evaluation errors in business or ceiling rules deny the affected action; unsupported constructs are rejected; hostile cost and cancellation limits stay bounded and request context stays isolated.}{yes}
\testrow{T-03}{R1-A, R1-B, R2}{Native first-hop authentication and rejection of adapter-asserted user identity; forged actor, subject or tenant; expired or widened delegation; audience mismatch; untrusted key locations; recipient substitution; signed-decision bytes, audience, expiry and replay checks; concurrent single-use consumption across replicas; signing-key rotation overlap. R1-B: recipient-key enrollment proof, fingerprint pinning and recipient substitution at publication and delivery (\rr{RCP-001}, \rr{RCP-002}, \rr{RCP-003}, \rr{RES-009}).}{yes}
\testrow{T-04}{R1-B, R2}{R1-B: changing a prepared storage operation, target version or approved protected-metadata change invalidates its grant; unsupported obligations fail before disclosure or commit; approvers see the canonical plan. R2 extends the same harness to tool arguments, destinations, models, amounts, redirects and registry changes.}{yes}
\testrow{T-05}{R1-B, R2}{R1-B: concurrent consumption of an approval for a storage operation permits at most its authorized use. R2 external-effect cases add budget reservation, dispatch ownership and reconciliation: UNKNOWN retains its reservation; ABANDONED retains the charge, consumed approval and redispatch prohibition.}{yes}
\testrow{T-06}{R1-A, R1-B, R1-C, R2}{One canonical-resource harness. R1-A: registry resolution, alias binding and compound checks. R1-B: native encrypted-object bindings and version races. R1-C: authorized unit ciphertext and envelope selection, denied-unit raw download and incomplete manifests. R2: MCP tools and resources preserve the same decisions and encrypted delivery boundaries, with decryption confined to the authorized client or local bridge.}{yes}
\testrow{T-07}{R3}{SigV4 canonicalization, signature or body mismatch, path encoding, duplicate keys and headers, skew, unsupported payload modes and scope mismatch behave as specified; named S3 clients pass with their documented settings.}{no}
\testrow{T-08}{R1-B, R1-C}{R1-B: object creation, version publication, concurrent compare-and-swap, conditional writes, deletion and garbage collection over the qualified backend. R1-C extends the same publication harness to staged units, interrupted manifest commits and retained document versions sharing a unit with distinct envelope contexts. No LIVE document references an unavailable or retired unit.}{yes}
\testrow{T-09}{R1-B, R1-C, R2}{R1-B: native object listing, stat, hidden-name and permitted metadata boundaries. R1-C: unit selection (plaintext in server mode, ciphertext and envelopes in client mode) and authenticated label and lineage metadata; client-mode parsing, snippets and search do not upload content plaintext, and server-mode ingestion keeps no plaintext index or snippet. R2: MCP returns only authorized content: server-mode plaintext, or client-mode ciphertext and envelopes. A model submission occurs only after client-side decryption and a separate authorized client disclosure; the server cannot verify arbitrary recipient behavior.}{yes}
\testrow{T-10}{R1-B, R1-C, R2}{Client mode. One client encryption/envelope harness covers wrong tenant, object/version, recipient and exact context bytes; corrupted, reordered or truncated ciphertext; and authenticated envelope provenance. The client encrypts before upload and decrypts after download; inspection after TLS termination at the platform shows only content ciphertext and opaque envelopes for client-mode content. R1-C adds unit/manifest bindings; R2 adds MCP client interoperability.}{no}
\testrow{T-11}{R1-B, R1-C, R2}{The broker delivers an existing opaque envelope only for the authorized recipient and exact object/version, with every signed prerequisite of Table~\ref{tab:keyops}. Missing envelopes, revoked recipients, pending or denied objects, stale decisions and broker/decision failure deny. Initial upload applies the same sharing prerequisites as later publication; target-recipient eligibility cannot be replaced by the publisher's rights. An envelope covers exactly one object/unit keyset, while a batch contains separate authorized envelopes. Server credentials cannot unwrap any client-mode data key. R1-C adds independent unit envelopes and prohibits obtaining a document-wide key from unit rights; R2 reuses the same broker cases through MCP.}{yes}
\testrow{T-12}{R1-B, R1-C}{Client mode. Authorized client key rotation and envelope replacement preserve intended reads. The server cannot create an envelope for a new recipient without an authorized key-holding client. Rewrap and revocation do not revoke already retained ciphertext and usable keys; lost client keys without independent recovery material remain unrecoverable. R1-C adds the same boundary for unit keys. Historical content-signing trust survives ordinary rotation; compromise and restore follow the explicit trust policy.}{yes}
\testrow{T-13}{R1-A, R1-B, R1-C, R2}{R1-A: admission and commit checkpoints honor expiry and authority revisions. R1-B: ciphertext streams and envelope delivery stop on revoked authority or overdue checks. R1-C: unit label-only changes with unchanged content generation. R2: managed-client execution-context revisions where the cooperating-client profile applies. Retained plaintext, envelopes, keys and already-emitted bytes keep the documented limits.}{yes}
\testrow{T-14}{R1-A}{A failed bundle activation does not publish readiness. Concurrent requests and batches keep one bundle. Stale nodes and invalid production configuration reject. Readiness cannot substitute for fresh admission. Rollback is a new higher epoch.}{yes}
\testrow{T-15}{R1-A, R1-B, R2, R3}{R1-A: mandatory audit records survive exporter failure and load. R1-B: failure of one qualified domain during ciphertext/envelope publication, envelope delivery or policy activation preserves acknowledged catalog, object and audit outcomes. R2 adds optional external tool dispatch; R3 adds custom-replica failures. Extend the same fault harness only at distinct durability boundaries.}{no}
\testrow{T-16}{R1-B, R2, R3}{R1-B: coordination and backend failures preserve primary fencing, quorum and revoked authority. A cancelled synchronous wait may expose an unreplicated row; no publication is acknowledged or envelope delivered against that unproven state until its durability barrier is established. R2 checks the barrier before optional external dispatch; R3 qualifies custom-storage partitions without quorum downgrade. The standby, synchronous-commit and failover cases need the \code{ha} profile; the restore case (\rr{AVL-008}) applies to both profiles.}{yes}
\testrow{T-17}{R3}{One shared repair and movement scenario covers corruption, lost wakeups, interrupted scanning and changed membership. Durable work resumes and replica count recovers without premature deletion.}{no}
\testrow{T-18}{R1-B, R1-C, R3}{Restore ciphertext, opaque envelopes, permitted catalog metadata and policies to a consistent point; quarantine inconsistent references and measure RPO/RTO. Server backups and temporary storage contain no recipient private keys, plaintext data keys or content plaintext; server-mode data keys appear only wrapped, and KEKs are recovered separately (\rr{SSE-006}). Test client-mode recovery separately using organization-controlled escrow or recovery recipients, never a server-held decryption credential. R1-C adds unit/manifest recovery; R3 exercises the custom backend. Missing client recovery material has an explicit unrecoverable outcome.}{no}
\testrow{T-19}{R1-A, R1-B}{Logs, traces, metrics and support bundles contain no secrets, raw keysets, inappropriate full arguments or cross-tenant identifiers. R1-B: client SDK diagnostics, envelope-service logs and catalog dumps contain no keysets, private keys or plaintext (\rr{KEY-007}, \rr{DAT-006}).}{no}
\testrow{T-20}{R1-A, R1-B, R1-C, R2, R3}{Role configuration, credentials and import allowlists reject a content-decryption or data-key unwrap path outside the storage role's and \code{keys} role's server-mode functions, any path that applies them to client-mode content, recipient private-key configuration and implicit public/anonymous grants. R1-B/R1-C qualify round trips in both modes, deny cases, failover and restore with MCP/models/tools absent; the server-side capture after TLS termination contains no client-mode plaintext. R2 checks MCP/client-bridge separation and independent scaling; R3 adds maintenance fences. Reuse the boundary harness.}{yes}
\testrow{T-21}{R1-C, R2}{R1-C: verify provenance and policy authorization of submitted labels and lineage records; the server does not claim to infer their truth from ciphertext. R2 uses a cooperating managed-client profile for conservative lineage and monotonic context marks, concurrent disclosure/egress, continuation and label-order changes. Test refusal of stale marks within that profile and document that arbitrary recipients can retain, relabel or disclose their decrypted content outside mediation.}{yes}
\testrow{T-22}{R1-A, R1-B, R1-C, R2}{Qualification of Section~\ref{sec:perf}. R1-A: decision limits and batching. R1-B: server-mode encryption/decryption throughput and key wrap/release rate, client encryption/decryption throughput and memory, envelope creation/delivery, ciphertext transfer, concurrent publication, listing, leases and degraded operation. R1-C: unit fan-out and recipient-envelope cardinality. R2: MCP and local-bridge overhead separately from storage; optional client model submission and tool latency are reported separately.}{no}
\testrow{T-23}{R1-B, R1-C, R2}{Encryption modes and server-side keys. R1-B: bucket mode at creation and change, policy visibility of the mode and a ceiling that requires one; refusal of cross-mode routes and of client material on server-mode versions; server-mode round trips with AAD substitution, corrupted, reordered and truncated ciphertext; key release only under a fresh exact decision (stale epoch, advanced head, other version, missing approval, revoked grant); KEK rotation without ciphertext change; capture after TLS termination shows plaintext only on authorized server-mode reads; no data key or plaintext in logs, catalog dumps, temporary files or backups; restore with missing KEK versions. R1-C: server-side splitting, unit-level release and complete-content release of originals. R2: MCP returns server-mode plaintext only after the same checks and client-mode content only as encrypted packages.}{yes}
\testrow{T-24}{R1-B, R1-C}{Share grants. Recording and removing grants requires \code{share}; grants outside the read set, beyond the grantor's delegation or to unknown or other-tenant principals reject; a grant without an honoring rule changes nothing; expiry and removal invalidate decisions and abort streams; ceilings and clearances still deny; in client mode a grant makes the grantee eligible but supplies no key. R1-C: grants on documents keep per-unit checks.}{yes}
\bottomrule
\end{longtable}

\subsection{Performance qualification}\label{sec:perf}
\reqn{PERF-001}{R1-B}{T-22}{Benchmarks measure decisions and context lookups separately from ciphertext transfer, server-mode encryption/decryption and key wrap/release, client encryption/decryption, recipient-envelope construction/delivery and optional tool latency. Server-mode key operations and payload cryptography are reported as their own baseline. They report hardware, policy and context sizes, policy selectivity, tenant skew, concurrency, payload sizes, cache configuration and failure mode, and include 1-, 10- and 50-resource batches, candidate-index selectivity, concurrent activation and preparation memory, cancellation, sparse visible listings and large streaming objects. They report p50, p95 and p99 and queue time, not only averages. Storage capacity tests include one failed qualified failure domain and backend recovery; custom-replica repair is included when R3 is enabled. Optional MCP, tool and model paths are measured separately and cannot stand in for native storage measurements. The initial engineering objectives are p95 under 10 ms for local pure policy evaluation and bounded memory per concurrent stream; end-to-end objectives are set after representative client, context and network measurements. These are targets for qualification, not existing capabilities, and any changed limit is recorded with why it preserves denial behavior and resource isolation.}

\subsection{Release gates}
\reqn{GATE-001}{P0}{REV}{A release is accepted only when its enabled surfaces meet the quality, provenance, dependency and enforcement-coverage requirements, and every applicable exit invariant has passing evidence for that release configuration. Case and test-family applicability uses minimum phase and enabled-feature prerequisites; disabled optional features do not gate unrelated surfaces. Scenario-only cases and missing, stale or failing evidence block qualification; document compilation and structural checks do not waive those blocks. The release includes contracts, build and dependency manifests, required notices, signed policy fixtures, operational guidance, and tested rollback and recovery procedures. The \code{--qualify} checker mode enforces the supplied case/evidence contract; independent review assesses the evidence itself and remaining release obligations.}
\reqn{GATE-002}{P0}{REV}{A Phase 0 decision (Section~\ref{sec:decisions}) is closed and recorded before the work it gates starts. Closing one does not reopen the core design: independent implementation, explicit delegation ceilings, one shared policy evaluation, mandatory adapters, immutable ciphertext, mode-bound content decryption (client-only in client mode, decision-bound key release in server mode), policy-controlled opaque-envelope delivery and fail-closed semantics remain required.}
\reqn{GATE-003}{P0}{REV}{R1-A, R1-B and R1-C are subphases of one release. No subphase is deployed to production by itself, and production use of R1 requires every R1 exit test to pass.}

\reqn{GATE-004}{R1-C}{T-20}{The R1 foundation passes its encrypted object and document round trips in both encryption modes, policy, share-grant, client/server key separation, isolation, durability, restore, independent client recovery and sizing gates with MCP, model-provider and external-tool services absent. MCP implementation begins only after this foundation is accepted. R2 qualifies the separate MCP service against server-mode plaintext reads and the client-mode ciphertext and opaque-envelope contract, with a compatible client SDK or local bridge for client-mode decryption. Model submission runs on an authorized client after it reads or decrypts content, not on a platform server; it and external tool execution are separately enabled extensions whose requirements and relevant cases must pass before use. Their absence does not block the ordinary MCP content service, and an unimplemented extension stays disabled.}
